\begin{table}[tbp]\centering
  \begin{threeparttable}
    \caption{Survey answers by the respondent's Chess.com rating}
    \label{tab:survey_by_rating}
    \begin{tabular}{lccc}
      \toprule
       & \multicolumn{1}{c}{Below median} & \multicolumn{1}{c}{Median or above} & \multicolumn{1}{c}{All}\\
      \midrule
      \textit{More helpful to review} & & & \\
      \quad A loss & 5 (38) & 11 (79) & 20 (57)\\
      \quad A win & 0 (0) & 0 (0) & 0 (0)\\
      \quad Both equally & 8 (62) & 2 (14) & 14 (40)\\
      \quad Neither & 0 (0) & 1 (7) & 1 (3)\\
      \textit{More pleasant to review} & & & \\
      \quad A win & 9 (69) & 10 (71) & 23 (66)\\
      \quad A loss & 3 (23) & 1 (7) & 6 (17)\\
      \quad No difference & 1 (8) & 3 (21) & 6 (17)\\
      \textit{Wish to shift review time} & & & \\
      \quad Toward losses & 8 (62) & 3 (21) & 13 (37)\\
      \quad Toward wins & 0 (0) & 0 (0) & 1 (3)\\
      \quad Same change for both & 5 (38) & 11 (79) & 21 (60)\\
       \\
      Respondents & 13 & 14 & 35\\
      \bottomrule
    \end{tabular}
    \begin{tablenotes}
      \footnotesize
      \item \textit{Note:} Counts, with percentages of the column in parentheses, by the respondent's rating relative to the median. The rating is the respondent's self-reported Chess.com rating; 27 of the 35 respondents reported one (median 1,404), the others entered 0, 1 or 3, which we treat as missing. Fisher's exact test of ``a loss'' against the other answers to the first question, by rating group: $p=0.05$.
    \end{tablenotes}
  \end{threeparttable}
\end{table}
